\section{Conclusion}\label{sec:conclusion}

Over four months the programme made elliptic-curve index calculus much
cheaper against itself.
\begin{itemize}
\item The decomposition oracle is two to five orders of magnitude cheaper
  where root finding applies.
\item The Gr\"obner engines are one to two orders cheaper in counted work,
  and the $\FF_p$ engine matches msolve.
\item The $n=83$ cold pipeline is $3.24\times$ faster after its set-up was
  stacked.
\item Two extension-field families moved their exponents: genus three
  after cubic GHS from $q^{1.91}$ to a local $q^{1.36}$, and $k=4$ with
  $D_4$ symmetry to a gap closing as $n^{-1/4}$.
\end{itemize}
By the success criterion each of these is a success.

Against rho, the honest summary is shorter. Cold, with a matched reference,
rho is cheaper on every elliptic curve the programme measured:
\begin{itemize}
\item $4.67$--$18.07\times$ on Koblitz curves of 37--47 bits;
\item $10.5$--$102\times$ on prime curves of 20--40 bits;
\item $55\times$ at Koblitz $n=83$.
\end{itemize}
On the families that are not deployed shapes, the gap is $1.15\times$
(genus three) and $4.3\times$ ($k=4$, one target). Every apparent win over
rho in the record disappeared once the reference was given equal advice,
priced alike, or matched in the same binary. Corrections of this kind
lengthened the measured distance more than any single improvement
shortened it.

Two things follow. First, the stage the programme improved most is no
longer the stage that dominates. Relation yield per probe and table set-up
now set the cost, and the next gains must come from there. Second, the
nearest approaches are on families where a fitted exponent is below rho's,
and both make a falsifiable prediction at sizes the harness can reach: a
$k=4$ crossover near $n\approx2^{68}$, and a genus-three crossover at
$q\le2^{22}$ once its rho is measured. Testing those two predictions, with
every row validated, is the work that would turn the ratchet into an
advance.
